%----------------------------------------------------------------------
% 中文摘要
%----------------------------------------------------------------------

% 把中文摘要寫在裡面
\begin{zhAbstract}
光的自旋角動量於 1936 年\add{首次於Beth的實驗中被量測}，而結構光所攜帶的軌道角動量則於 1992 年\add{由Allen等人}透過拉蓋爾–高斯(Laguerre–Gaussian)光束\add{於實驗中}驗證。近年研究指出，在非近軸條件下，光場中可能出現自旋與軌道角動量之間的耦合。然而，僅在近軸近似下，光束才能同時具備明確的自旋與軌道角動量。\add{最近的}研究已指出，

\delete{自旋–軌道交互作用出現於近軸近似下蓋爾–高斯光束中微弱但非零的縱向分量之相位結構中。本研究進一步透過疊加兩道蓋爾–高斯光束形成向量渦旋光(vector vortex beams)，並證明其對應的自旋–軌道交互作用項可延續並保留於縱向場分量的振幅之中。}
\add{增加SOC重要性}

\delete{為了提供觀測此自旋–軌道交互作用特徵的可行途徑，二維過渡金屬二硫族化物(transition metal dichalcogenides)成為研究光與物質交互作用的理想平台之一，這主要歸因於其在可見光至近紅外波段具有直接能隙以及強烈的激子效應。其中，灰激子(gray excitons)因可透過光場的縱向分量與光耦合，特別適合用於探測此 向量渦旋光中自旋–軌道交互作用的特徵。此外，灰激子的躍遷偶極矩在動量空間中近似於均勻分布，使得光與物質的交互作用更能直接反映向量渦旋光所蘊含的自旋–軌道交互作用特徵。本研究發現，谷混合亮激子（valley–mixed bright excitons）與灰激子（gray excitons）皆可透過不同機制，利用向量渦旋光的相對相位調控，實現選擇性激發。}

\add{最後}，本研究延伸先前對谷混合激子波包的探討，\add{分析}激子內部波函數中的相位結構，使得任意相位無法被允許提出並影響波包形貌。


\delete{雖然對於谷激子波包，相位結構在一定近似下仍可分離，但對於谷混合激子而言，問題更為複雜。本研究發現在此情況下，除了要寫下並提取相位結構，更需近一步建構激子波函數的內部結構，才能避免任意相位因子對最終波包形貌造成顯著影響。}
\add{激子相位對拓樸的重要性}
    % 這個Command會自動幫你把Config裡面設定的東東填進來
    \zhAbsKeywords
\end{zhAbstract}
